% Suggested revision in response to the reviewer comment on calibration.
% This file is intentionally separate from body.tex.  The text below can be
% adapted after the experimental pulsed-NMR and CPMG analysis is available.

% -----------------------------------------------------------------------------
% Suggested response to the reviewer
% -----------------------------------------------------------------------------
% Thank you for this helpful suggestion.  We have revised the calibration
% section to make the expected behavior and quantitative agreement more
% explicit.  For each calibration, the simulated signal is now compared with
% the corresponding analytical expectation, and the residual (simulation
% minus theory) is shown in an additional row of the relevant figure.  We
% also added an appendix describing a common analysis workflow for comparing
% simulated signals with measured pulsed-NMR and CPMG data.  The comparison
% includes the same pulse sequence, sampling interval, receiver phase,
% normalization, and relaxation parameters in simulation and experiment.
% Experimental-data comparisons will be populated with the measured data and
% their uncertainties; where appropriate, agreement is quantified using a
% normalized residual or reduced chi-square.

% -----------------------------------------------------------------------------
% Suggested replacement/addition for the calibration section
% -----------------------------------------------------------------------------

% Replace or substantially expand the current opening paragraphs of the
% Calibration section with the following text.

\subsection{Calibration}
\label{sec:calibration_revision}

Calibration was performed at two levels.  First, idealized pulse sequences
and continuous-wave excitation were compared with analytical solutions of
the Bloch equations.  This tests the numerical integrator and the mapping
between input parameters and observables.  Second, the simulated observables
were prepared in the same form as experimental NMR data, providing a direct
workflow for comparison with measured free-induction decays and spin echoes.
For every calibration, the simulated and expected signals were evaluated on
the same time grid.  In addition to displaying the signals themselves, we
show the pointwise residual
\begin{equation}
    r(t_i) = M_{\mathrm{simu}}(t_i)-M_{\mathrm{theory}}(t_i),
    \label{eq:calibration_residual_revision}
\end{equation}
and summarize its normalized squared norm,
\begin{equation}
    \chi^2_{\mathrm{rel}}
    = \frac{\sum_i |r(t_i)|^2}
    {\sum_i |M_{\mathrm{theory}}(t_i)|^2}.
    \label{eq:calibration_chisq_revision}
\end{equation}
The residual panels make numerical errors, pulse-discretization effects, and
any systematic deviations from the analytical model visible rather than
representing them only by a single summary statistic.

For the pulsed-NMR calibration, the analytical expectations used for the
comparison are the free-induction decay
\begin{equation}
    M_{\perp}(t) = A_{\mathrm{FID}}
    \exp(-t/T_2^*)\exp\!\left[i(\omega_{\mathrm{L}}t+\phi)\right],
    \label{eq:fid_expectation_revision}
\end{equation}
and the spin-echo envelope
\begin{equation}
    A_{\mathrm{echo}}(t_{\mathrm{echo}})
    = A_0\exp(-t_{\mathrm{echo}}/T_2).
    \label{eq:echo_expectation_revision}
\end{equation}
Here, $\omega_{\mathrm{L}}=\gamma B_0$ and the phase $\phi$ is chosen to
match the receiver reference.  For the weak, resonant continuous-wave case,
the transverse response is compared with
\begin{equation}
    M_{\perp}(t) \simeq \gamma B_1 T_2^*
    \left(1-\exp(-t/T_2^*)\right),
    \qquad \gamma B_1T_2^*\ll 1.
    \label{eq:cw_expectation_revision}
\end{equation}
For longitudinal relaxation of a hyperpolarized sample, the expected curve is
\begin{equation}
    M_z(t)=M_{\mathrm{eq}}
    +\left[M_z(0)-M_{\mathrm{eq}}\right]\exp(-t/T_1).
    \label{eq:t1_expectation_revision}
\end{equation}

The calibration parameter ranges and numerical results are listed in
Table~\ref{tab:calibration_summary_revision}.  The integration step size is
reported explicitly because the residual is expected to decrease as the step
size is reduced, while the computational cost increases.

% -----------------------------------------------------------------------------
% Suggested figure organization
% -----------------------------------------------------------------------------

% Regenerate the existing calibration figures with three rows where practical:
% (i) applied field/pulse sequence, (ii) simulated and theoretical signal,
% and (iii) residual.  For example, replace the existing SpinEcho_and_CW.pdf
% figure with a version containing the following caption.

\begin{figure}[htb]
    \centering
    \includegraphics[width=0.99\linewidth]{figures/SpinEcho_and_CW_with_residuals.pdf}
    \caption{Calibration against analytical expectations.  The upper rows
    show the applied fields and the middle rows show the simulated
    magnetization together with the corresponding theoretical curves.  The
    bottom rows show the pointwise residual defined in
    Eq.~\eqref{eq:calibration_residual_revision}.  Panels (a) and (b) show,
    respectively, the pulsed-NMR free-induction/spin-echo sequence and the
    weak resonant continuous-wave response.  Signal amplitudes are normalized
    consistently between the simulation and analytical model.}
    \label{fig:calibration_with_residuals_revision}
\end{figure}

% For the T1 calibration, use the same layout with the simulated M_z curve,
% exponential expectation, and residual as the three rows.

\begin{figure}[htb]
    \centering
    \includegraphics[width=0.70\linewidth]{figures/T1_hyperpolarized_with_residual.pdf}
    \caption{Longitudinal relaxation calibration.  The simulated $M_z$
    signal is compared with the exponential expectation in
    Eq.~\eqref{eq:t1_expectation_revision}; the lower panel shows their
    residual.  The magnetization is normalized by $M_{\mathrm{eq}}$.}
    \label{fig:calibration_T1_residual_revision}
\end{figure}

% Add a residual column to the calibration summary table if the numerical
% values are available from the revised scripts.

\begin{table}[htb]
\small
\caption{Expanded summary of the numerical calibration tests.  The relative
$\chi^2_{\mathrm{rel}}$ is defined in Eq.~\eqref{eq:calibration_chisq_revision}.
The entries marked ``TBD'' should be replaced by values generated using the
same parameter sets as the revised figures.}
\label{tab:calibration_summary_revision}
\resizebox{\linewidth}{!}{%
\begin{tabular}{llllll}
    \toprule
    \textbf{Calibration} & \textbf{Observable} &
    \textbf{Theory used} & \textbf{Parameter range} &
    $\bm{\chi^2_{\mathrm{rel}}}$ & \textbf{Main limitation} \\
    \midrule
    Pulsed NMR: FID & $M_{\perp}(t)$ &
    Eq.~\eqref{eq:fid_expectation_revision} &
    $T_2^*$, $\nu_{\mathrm L}$ & TBD & Pulse discretization \\
    Pulsed NMR: spin echo & Echo envelope &
    Eq.~\eqref{eq:echo_expectation_revision} &
    $T_2$, pulse spacing & TBD & Finite pulse width \\
    CW NMR & $M_{\perp}(t)$ &
    Eq.~\eqref{eq:cw_expectation_revision} &
    $B_1$, $T_2^*$ & TBD & Weak-drive approximation \\
    Hyperpolarized sample & $M_z(t)$ &
    Eq.~\eqref{eq:t1_expectation_revision} &
    $T_1$, $M_z(0)$ & TBD & Step-size error \\
    \bottomrule
\end{tabular}}
\end{table}

% -----------------------------------------------------------------------------
% Proposed appendix for comparison with real measurements
% -----------------------------------------------------------------------------

\appendix
\section{Comparison with measured pulsed-NMR and CPMG signals}
\label{app:experimental_comparison_revision}

This appendix describes the procedure used to compare \axionbloch{} outputs
with measured NMR data.  The purpose is to separate discrepancies caused by
the numerical solver from those caused by experimental imperfections or by
an incomplete physical model.

\subsection{Data and metadata}

Each data set should be accompanied by the static-field magnitude and
orientation, gyromagnetic ratio, pulse amplitudes and phases, pulse duration,
interpulse delay, sampling interval, receiver phase, number of averages, and
the values or independent estimates of $T_1$, $T_2$, and $T_2^*$.  The exact
pulse sequence is reproduced in the simulation, including finite pulse
duration when it is relevant.  The measured and simulated traces are then
resampled onto a common time grid.

\subsection{Pulsed-NMR comparison}

For a free-induction decay, the comparison should use the same acquisition
window and account for receiver phase, dead time, baseline offset, and signal
normalization.  The fitted Larmor frequency and $T_2^*$ can first be used to
check the phase evolution and envelope independently.  For CPMG data, the
echo centers and echo amplitudes should be extracted using an identical
procedure for simulation and measurement.  The echo amplitudes can then be
compared with the $T_2$ envelope in Eq.~\eqref{eq:echo_expectation_revision}.

The comparison should be shown in three panels: measured and simulated
signals, the residual, and (where uncertainties are available) the residual
divided by the experimental standard deviation.  A suitable weighted metric
is
\begin{equation}
    \chi^2_{\nu} = \frac{1}{N-p}
    \sum_{i=1}^{N}
    \frac{\left[M_{\mathrm{meas}}(t_i)-M_{\mathrm{simu}}(t_i)\right]^2}
    {\sigma_i^2},
    \label{eq:experimental_reduced_chisq_revision}
\end{equation}
where $N$ is the number of data points, $p$ is the number of fitted
parameters, and $\sigma_i$ is the uncertainty of point $i$.  If pointwise
uncertainties are unavailable, the unweighted residual norm in
Eq.~\eqref{eq:calibration_chisq_revision} should be reported instead, with
that limitation stated explicitly.

\subsection{Scope and limitations}

The first experimental comparison should use a calibration sample and a
well-characterized pulse sequence.  Agreement is expected to be limited by
field inhomogeneity, pulse imperfections, receiver response, noise, finite
bandwidth, and possible deviations from single-component Bloch relaxation.
These effects should not be absorbed into the numerical-calibration metric
without being identified.  The appendix therefore reports both the ideal
analytical comparison and the experiment comparison, using separate residual
definitions and clearly stated fitted parameters.

% End of proposed revision.
